\setcounter{section}{27}

  \zwischenueberschrift{Soal latihan}

\inputaufgabe
{}
{

Misalkan
\maabb {p} { E } { M
} {}
suatu
\definitionsverweis {bundel vektor diferensiabel}{}{}
di atas
\definitionsverweis {manifold diferensiabel}{}{}
$M$, yang dilengkapi dengan
\definitionsverweis {koneksi linear}{}{}
$\nabla$. Misalkan $L$ manifold diferensiabel lain dan
\maabb {\psi} { L } { M
} {}
suatu
\definitionsverweis {pemetaan terdiferensialkan secara kontinu}{}{}
dengan
\definitionsverweis {koneksi tarikan balik}{}{}
$\psi^*\nabla$ pada
\definitionsverweis {bundel vektor tarikan balik}{}{}
$\psi^*E$ di atas $L$. Tunjukkan bahwa suatu penampang yang terdiferensialkan secara kontinu
\maabb {s} { L } { E
} {}
bersifat
\definitionsverweis {horizontal}{}{}
terhadap $\nabla$ jika dan hanya jika penampang terkait
\maabbdisp {s} { L } { \psi^*E
} {}
bersifat horizontal terhadap $\psi^* \nabla$.

}
{} {}

\inputaufgabe
{}
{

Misalkan
\maabb {p} {M \times \R^r } { M
} {}
adalah
\definitionsverweis {bundel vektor}{}{}
trivial di atas suatu
\definitionsverweis {manifold diferensiabel}{}{}
$M$, dilengkapi dengan
\definitionsverweis {koneksi trivial}{}{}
$\nabla$. Misalkan $L$ manifold diferensiabel lain dan
\maabb {\psi} { L } { M
} {}
suatu
\definitionsverweis {pemetaan terdiferensialkan secara kontinu}{}{.}
Tunjukkan bahwa
\definitionsverweis {koneksi tarikan balik}{}{}
$\psi^*\nabla$
\zusatzklammer {pada $L \times \R^r$} {} {}
juga trivial.

}
{} {}

\inputaufgabe
{}
{

Misalkan
\maabb {p} { E } { M
} {}
suatu
\definitionsverweis {bundel vektor diferensiabel}{}{}
di atas
\definitionsverweis {manifold diferensiabel}{}{}
$M$, yang dilengkapi dengan
\definitionsverweis {koneksi linear}{}{}
$\nabla$. Misalkan

\mavergleichskette
{\vergleichskette
{ I
}
{ \subseteq }{ \R
}
{  }{
}
{    }{
}
{   }{
}
}
{}{}{}
suatu
\definitionsverweis {interval terbuka}{}{}
dan
\maabb {\gamma} { I } { M
} {}
suatu
\definitionsverweis {kurva terdiferensialkan secara kontinu}{}{}
dengan
\definitionsverweis {koneksi tarikan balik}{}{}
$\gamma^*\nabla$ pada
\definitionsverweis {bundel vektor tarikan balik}{}{}
$\gamma^*E$ di atas $I$. Tentukan
\definitionsverweis {simbol-simbol Christoffel}{}{}
dari
\mathl{\gamma^* \nabla}{.}

}
{} {}

\inputaufgabegibtloesung
{}
{

Kita tinjau parametrisasi trigonometris
\maabbeledisp {\psi} {\R } { S^1
} { t } { \left(  \cos  t   , \,   \sin  t                   \right)
} {}
dari lingkaran satuan. Tetapkan

\mavergleichskette
{\vergleichskette
{ c
}
{ \in }{ [0, 2 \pi[
}
{  }{
}
{    }{
}
{   }{
}
}
{}{}{}
. Pada bundel vektor trivial
\maabbdisp {} {S^1 \times \R^2} { S^1
} {}
diberikan suatu
\definitionsverweis {koneksi linear}{}{}
$\nabla$ sedemikian sehingga sepanjang $\psi$
\definitionsverweis {koneksi tarikan balik}{}{}
$\psi^* \nabla$ diberikan oleh
\definitionsverweis {simbol-simbol Christoffel}{}{}
\zusatzklammer {indeks untuk satu-satunya arah penurunan tidak dituliskan} {} {}

\mavergleichskettedisp
{\vergleichskette
{ \begin{pmatrix} \tilde{\Gamma}_1^1 (t)   &  \tilde{\Gamma}_1^2 (t)   \\ \tilde{\Gamma}_2^1(t)  &  \tilde{\Gamma}_2^2(t)  \end{pmatrix}
}
{ =} { \begin{pmatrix}  0  &   - { \frac{  c  }{  2 \pi   } }   \\  { \frac{  c  }{  2 \pi  } }  &  0  \end{pmatrix}
}
{ } {
}
{ } {
}
{ } {
}
}
{}{}{}
. Untuk suatu titik

\mavergleichskette
{\vergleichskette
{ Q
}
{ = }{ \begin{pmatrix}  a  \\b   \end{pmatrix}
}
{ \in }{ \R^2
}
{    }{
}
{   }{
}
}
{}{}{}
kita tinjau pemetaan
\maabbeledisp {\Psi_{Q}} {\R} { S^1 \times \R^2
} { t } { \left(  \cos  t   , \,   \sin  t  , \,   M(t)  \begin{pmatrix}  a  \\b   \end{pmatrix}                 \right)
} {}
dengan

\mavergleichskettedisp
{\vergleichskette
{ M(t)
}
{ =} { \begin{pmatrix}  \cos  { \frac{  ct  }{  2 \pi   } }   &   -  \sin  { \frac{  ct  }{  2 \pi   } }    \\  \sin  { \frac{  ct  }{  2 \pi  } }   &  \cos  { \frac{  ct  }{  2 \pi   } }  \end{pmatrix}
}
{ } {
}
{ } {
}
{ } {
}
}
{}{}{.}
\aufzaehlungfuenf{Tentukan $\Psi_Q(0)$.
}{Tentukan $\Psi_Q(2 \pi)$.
}{Tunjukkan bahwa $\Psi_Q$ adalah suatu
\definitionsverweis {pengangkatan horizontal}{}{}
sepanjang $\psi$.
}{Tunjukkan bahwa
\maabbeledisp {} {\R^2} { \R^2
} { Q } { \Psi_Q(2 \pi)
} {,}
merupakan
\definitionsverweis {isometri linear}{}{}
\zusatzklammer {titik dasar $(1,0)$ tidak dituliskan di sini} {} {.}
}{Tunjukkan bahwa
\maabbeledisp {} {\Z} { \operatorname{SL}_{  2  } \! \left(  \R  \right)
} { k } { \left(  Q \mapsto \Psi_Q(2k \pi )  \right)
} {,}
merupakan
\definitionsverweis {homomorfisme grup}{}{.}
}

}
{} {}

\inputaufgabegibtloesung
{}
{

Untuk kurva
\maabbeledisp {\psi} {\R} { {\mathbb H}
} { t } { \left(  t   , \,  e^t                   \right)
} {,}
yang bernilai dalam
\definitionsverweis {setengah bidang Poincaré}{}{}
${\mathbb H}$, tentukan
\definitionsverweis {simbol-simbol Christoffel}{}{}
untuk
\definitionsverweis {koneksi tarikan balik}{}{}
\zusatzklammer {dari
\definitionsverweis {koneksi Levi-Civita}{}{}
pada ${\mathbb H}$} {} {}
di atas $\R$.

}
{} {}

Dalam soal berikut, perhatikan bahwa menurut
[[Topologische Räume/Relatives Produkt/Hintereinanderschaltung/Aufgabe|Soal 11.23]],
tidak ada bedanya apakah tarikan balik bundel vektor ke $Y$ ditentukan langsung atau melalui $Z$.

\inputaufgabe
{}
{

Misalkan
\maabb {p} { E } { M
} {}
suatu
\definitionsverweis {bundel vektor diferensiabel}{}{}
di atas
\definitionsverweis {manifold diferensiabel}{}{}
$M$, yang dilengkapi dengan
\definitionsverweis {koneksi linear}{}{}
$\nabla$. Misalkan $Y,Z$ manifold-manifold diferensiabel lain dengan
\definitionsverweis {pemetaan-pemetaan terdiferensialkan secara kontinu}{}{}

\mathdisp {Y \stackrel{\varphi}{\longrightarrow} Z \stackrel{\psi}{\longrightarrow}  M} { . }

Tunjukkan bahwa untuk
\definitionsverweis {koneksi-koneksi tarikan balik}{}{}
pada
\mathl{\left(  \psi \circ \varphi  \right)^* E}{} berlaku hubungan

\mavergleichskettedisp
{\vergleichskette
{ \left(  \psi \circ \varphi  \right)^*(\nabla)
}
{ =} { \varphi^*  \left(  \psi^*\nabla  \right)
}
{ } {
}
{ } {
}
{ } {
}
}
{}{}{.}

}
{} {}

\inputaufgabe
{}
{

Misalkan $M$ suatu
\definitionsverweis {manifold Riemann}{}{}
yang dilengkapi dengan
\definitionsverweis {koneksi Levi-Civita}{}{}
dan misalkan
\maabbdisp {\gamma} { I } { M
} {}
suatu
\definitionsverweis {kurva}{}{}
yang dua kali terdiferensialkan secara kontinu. Tunjukkan bahwa $\gamma$ adalah
\definitionsverweis {kurva geodesik}{}{}
jika dan hanya jika kurva turunannya
\maabb {\gamma'} { I } { TM
} {}
merupakan
\definitionsverweis {penampang horizontal}{}{}
dalam
\definitionsverweis {bundel tangen}{}{}
$TM$.

}
{} {}

\inputaufgabe
{}
{

Misalkan

\mavergleichskette
{\vergleichskette
{ I
}
{ \subseteq }{ \R
}
{  }{
}
{    }{
}
{   }{
}
}
{}{}{}
suatu
\definitionsverweis {interval terbuka}{}{}
dengan fungsi positif yang terdiferensialkan secara kontinu
\maabbdisp {g} { I } { \R_+
} {,}
yang kita pandang sebagai
\definitionsverweis {metrik Riemann}{}{}
pada $I$. Karakterisasikan
\definitionsverweis {kurva-kurva geodesik}{}{}
\maabbdisp {\gamma} { J } { I
} {,}
dengan $J$ suatu interval lain.

}
{} {}

\inputaufgabegibtloesung
{}
{

Untuk kurva
\maabbeledisp {\psi} {\R} { {\mathbb H}
} { t } { \left(  t   , \,  e^t                   \right)
} {,}
yang bernilai dalam
\definitionsverweis {setengah bidang Poincaré}{}{}
${\mathbb H}$, tentukan
\definitionsverweis {percepatan tangensial}{}{}
untuk setiap $t$.

}
{} {}

\inputaufgabe
{}
{

Berikan contoh sebuah
\definitionsverweis {manifold Riemann}{}{}
yang
\definitionsverweis {terhubung}{}{}
$M$ dan dua titik

\mavergleichskette
{\vergleichskette
{ P,Q
}
{ \in }{ M
}
{  }{
}
{    }{
}
{   }{
}
}
{}{}{,}
yang tidak dapat dihubungkan oleh suatu
\definitionsverweis {kurva geodesik}{}{.}

}
{} {}

\inputaufgabe
{}
{

Berikan contoh sebuah
\definitionsverweis {manifold Riemann}{}{}
yang
\definitionsverweis {terhubung}{}{}
$M$ dan dua titik

\mavergleichskette
{\vergleichskette
{ P,Q
}
{ \in }{ M
}
{  }{
}
{    }{
}
{   }{
}
}
{}{}{,}
yang dapat dihubungkan oleh beberapa
\definitionsverweis {kurva geodesik}{}{}
yang semuanya mempunyai panjang sama.

}
{} {}

\inputaufgabe
{}
{

Berikan contoh sebuah
\definitionsverweis {manifold Riemann}{}{}
yang
\definitionsverweis {terhubung}{}{}
$M$ dan dua titik

\mavergleichskette
{\vergleichskette
{ P,Q
}
{ \in }{ M
}
{  }{
}
{    }{
}
{   }{
}
}
{}{}{,}
yang dapat dihubungkan oleh beberapa
\definitionsverweis {kurva geodesik}{}{}
yang panjangnya tidak semuanya sama.

}
{} {}

\inputaufgabegibtloesung
{}
{

Verifikasikan bahwa sistem persamaan diferensial biasa orde dua

\mavergleichskettedisp
{\vergleichskette
{ \gamma_1^{\prime \prime} (t)
}
{ =} { { \frac{   2  }{  \gamma_2(t) } } \gamma_1'(t) \gamma_2'(t)
}
{ } {
}
{ } {
}
{ } {
}
}
{}{}{}
dan

\mavergleichskettedisp
{\vergleichskette
{ \gamma_2^{\prime \prime} (t)
}
{ =} { -  { \frac{   1  }{  \gamma_2(t) } }  \gamma_1'(t) \gamma_1'(t) + { \frac{   1  }{  \gamma_2(t) } }    \gamma_2'(t) \gamma_2'(t)
}
{ } {
}
{ } {
}
{ } {
}
}
{}{}{}
diselesaikan oleh kurva-kurva

\mavergleichskettedisp
{\vergleichskette
{ \gamma(t)
}
{ =} { \begin{pmatrix}  s+r \tanh  t  \\ { \frac{   r  }{  \cosh  t  } }    \end{pmatrix}
}
{ } {
}
{ } {
}
{ } {
}
}
{}{}{}
untuk

\mavergleichskette
{\vergleichskette
{ s
}
{ \in }{ \R
}
{  }{
}
{    }{
}
{   }{
}
}
{}{}{,}

\mavergleichskette
{\vergleichskette
{ r
}
{ \in }{ \R_+
}
{  }{
}
{    }{
}
{   }{
}
}
{}{}{.}

}
{} {}

\inputaufgabe
{}
{

Kita tinjau lingkaran-lingkaran yang berpusat di
\mathl{(s,0)}{} dan berjari-jari

\mavergleichskette
{\vergleichskette
{ r
}
{ > }{ 0
}
{  }{
}
{    }{
}
{   }{
}
}
{}{}{.}
Karakterisasikan kapan dua lingkaran semacam ini berpotongan.

}
{} {}

\inputaufgabe
{}
{

Misalkan

\mavergleichskette
{\vergleichskette
{ P
}
{ = }{(a,b)
}
{ \in }{ {\mathbb H}
}
{    }{
}
{   }{
}
}
{}{}{}
suatu titik dalam
\definitionsverweis {setengah bidang Poincaré}{}{}
dan misalkan

\mavergleichskette
{\vergleichskette
{ v
}
{ = }{ (v_1,v_2)
}
{ \in }{ \R^2
}
{    }{
}
{   }{
}
}
{}{}{}
suatu
\definitionsverweis {vektor tangen}{}{}
di $P$.
\aufzaehlungzwei {Tentukan sebuah setengah lingkaran berpusat pada sumbu $x$
\zusatzklammer {dengan garis-garis vertikal sebagai kasus ekstrem} {} {,}
yang melalui $P$ dan menyinggung $\R v$.
} {Tentukan suatu
\definitionsverweis {kurva geodesik}{}{,}
yang merealisasikan vektor tangen ini.
}

}
{} {}

\inputaufgabe
{}
{

Misalkan
\maabb {\psi} { L } { M
} {}
suatu
\definitionsverweis {isometri}{}{}
antara
\definitionsverweis {manifold-manifold Riemann}{}{}
\mathkor {} {L} {dan} {M} {.}
\aufzaehlungdrei{Tunjukkan bahwa
\definitionsverweis {koneksi tarikan balik}{}{}
dari
\definitionsverweis {koneksi Levi-Civita}{}{}
pada $M$ sama dengan koneksi Levi-Civita pada $L$.
}{Tunjukkan bahwa
\definitionsverweis {percepatan-percepatan tangensial}{}{}
dari kurva yang dua kali terdiferensialkan
\maabb {\gamma} { I } { L
} {}
dan, berturut-turut,
\maabb {\psi \circ \gamma} { I } { M
} {}
saling bersesuaian.
}{Tunjukkan bahwa
\definitionsverweis {kurva-kurva geodesik}{}{}
di $L$ dan di $M$ saling bersesuaian.
}

}
{} {}

Misalkan $M$ suatu
\definitionsverweis {manifold Riemann}{}{,}
dengan
\definitionsverweis {bundel tangen}{}{}
$TM$ yang dilengkapi
\definitionsverweis {koneksi Levi-Civita}{}{.}
Misalkan
\maabb {\gamma} { I } { M
} {}
suatu
\definitionsverweis {kurva diferensiabel}{}{.}
Suatu objek diferensiabel yang didefinisikan sepanjang $\gamma$ dan merupakan
\definitionsverweis {medan vektor}{}{},
\maabb {F} { I } { TM
} {}
disebut
\definitionswort {paralel sepanjang}{}
$\gamma$ jika

\mavergleichskettedisp
{\vergleichskette
{ (\nabla_{\partial} F)(t)
}
{ =} { 0
}
{ } {
}
{ } {
}
{ } {
}
}
{}{}{}
untuk semua

\mavergleichskette
{\vergleichskette
{ t
}
{ \in }{ I
}
{  }{
}
{    }{
}
{   }{
}
}
{}{}{.}

Untuk soal berikut, bandingkan dengan
[[Differenzierbare Hyperfläche/Zusammenhang/Kurve/Parallel und horizontal/Aufgabe|Soal 25.13]].

\inputaufgabe
{}
{

Misalkan

\mavergleichskette
{\vergleichskette
{ Y
}
{ \subseteq }{ W
}
{ \subseteq }{ \R^n
}
{    }{
}
{   }{
}
}
{}{}{,}
dengan $W$
\definitionsverweis {terbuka}{}{,}
dan suatu
\definitionsverweis {hiperpermukaan diferensiabel}{}{}
yang juga kita pandang sebagai
\definitionsverweis {manifold Riemann}{}{}
melalui ruang sekitar $\R^n$ dengan
\definitionsverweis {hasil kali skalar standar}{}{.}
Misalkan
\maabb {\gamma} { I } { Y
} {}
suatu
\definitionsverweis {kurva diferensiabel}{}{}
dan
\maabb {F} { I } { TY \subseteq Y \times \R^n
} {}
suatu
\definitionsverweis {medan vektor}{}{}
diferensiabel sepanjang $\gamma$. Tunjukkan bahwa $F$
\definitionswort {paralel sepanjang}{}
$\gamma$ dalam arti
[[Hyperfläche/Kurve/Paralleles Vektorfeld/Definition|Definisi 6.6]]
jika dan hanya jika $F$
\definitionsverweis {paralel}{}{}
sepanjang $\gamma$ terhadap
\definitionsverweis {koneksi Levi-Civita}{}{}
pada $TY$.

}
{} {}

\inputaufgabe
{}
{

Misalkan $M$ suatu
\definitionsverweis {manifold Riemann}{}{,}
dengan
\definitionsverweis {bundel tangen}{}{}
$TM$ yang dilengkapi
\definitionsverweis {koneksi Levi-Civita}{}{.}
Misalkan
\maabb {\gamma} { I } { M
} {}
suatu kurva yang dua kali
\definitionsverweis {terdiferensialkan}{}{.}
Tunjukkan bahwa $\gamma$ adalah
\definitionsverweis {kurva geodesik}{}{}
jika dan hanya jika turunannya
\maabbdisp {\gamma'} { I } { TM
} {}
merupakan
\definitionsverweis {medan vektor paralel}{}{.}

}
{} {}

\inputaufgabe
{}
{

Tunjukkan bahwa pada
\definitionsverweis {ruang Euklides}{}{}
$\R^n$, 
\definitionsverweis {metrik Riemann}{}{}
sama dengan
\definitionsverweis {jarak Euklides}{}{.}

}
{} {}

  \zwischenueberschrift{Tugas yang dikumpulkan}

\inputaufgabe
{4}
{

Kita tinjau $\R_+$ dengan struktur yang diberikan oleh

\mavergleichskette
{\vergleichskette
{ g(t)
}
{ = }{ t^2
}
{  }{
}
{    }{
}
{   }{
}
}
{}{}{,}
yaitu suatu
\definitionsverweis {metrik Riemann}{}{.}
Tentukan
\definitionsverweis {kurva-kurva geodesik}{}{}
di $\R_+$.

}
{} {}

\inputaufgabe
{6 (3+3)}
{

Kita tinjau kurva
\maabbeledisp {\psi} {\R} { {\mathbb H}
} { t } { \left(  t   , \,   1+t^2                   \right)
} {,}
yang bernilai dalam
\definitionsverweis {setengah bidang Poincaré}{}{}
${\mathbb H}$.
\aufzaehlungzwei {Tentukan
\definitionsverweis {simbol-simbol Christoffel}{}{}
untuk
\definitionsverweis {koneksi tarikan balik}{}{}
\zusatzklammer {dari
\definitionsverweis {koneksi Levi-Civita}{}{}
pada ${\mathbb H}$} {} {}
di atas $\R$.
} {Tentukan
\definitionsverweis {percepatan tangensial}{}{}
dari $\psi$ untuk setiap $t$.
}

}
{} {}

\inputaufgabe
{4}
{

Tentukan
\definitionsverweis {kurva-kurva geodesik}{}{}
pada
\definitionsverweis {cakram hiperbolik}{}{.}

}
{} {}

\inputaufgabe
{4}
{

Misalkan

\mavergleichskette
{\vergleichskette
{ n
}
{ \geq }{ 2
}
{  }{
}
{    }{
}
{   }{
}
}
{}{}{}
dan

\mavergleichskettedisp
{\vergleichskette
{ M
}
{ =} { \R^n \setminus \{P_1  , \ldots , P_m \}
}
{ \subseteq} { \R^n
}
{ } {
}
{ } {
}
}
{}{}{,}
yang dipandang sebagai
\definitionsverweis {manifold Riemann}{}{}
dengan
\definitionsverweis {hasil kali skalar standar}{}{}
terinduksi. Tunjukkan bahwa pada $M$,
\definitionsverweis {metrik Riemann}{}{}
sama dengan
\definitionsverweis {jarak Euklides}{}{.}

}
{} {}

 [[Kategorie:Latexseite]]
